\documentclass[11pt,a4paper]{article}
\usepackage[utf8]{inputenc}
\usepackage[T1]{fontenc}
\usepackage{amsmath,amsfonts,amssymb}
\usepackage{graphicx}
\usepackage{hyperref}
\usepackage{listings}
\usepackage{color}
\usepackage{booktabs}
\usepackage{float}
\usepackage{geometry}
\geometry{margin=1in}
\definecolor{zenblue}{RGB}{41,121,255}
\definecolor{zengreen}{RGB}{52,199,89}
\hypersetup{colorlinks=true,linkcolor=zenblue,urlcolor=zenblue,citecolor=zenblue}

\title{\textbf{Zen: An Instruction-Following Language Model Built on Qwen3-8B}\\
\large Technical Report v2025.01}
\author{Antje Worring, Zach Kelling \\ Zen LM Research Team\\
\texttt{research@zenlm.org}}
\date{January 2025}

\begin{document}
\maketitle

\begin{abstract}
We present \textbf{Zen}, an 8-billion-parameter instruction-following language model that serves as
the foundation of the Zen family. Zen is \emph{not} trained from scratch: it is a fine-tuned and
repackaged derivative of \textbf{Qwen3-8B} \cite{qwen3report}, the open-weight dense model released
by Alibaba's Qwen team under the Apache-2.0 license. Our contribution is post-training and
packaging---supervised fine-tuning (SFT) and preference optimization on curated instruction data,
together with deployment artifacts (GGUF, MLX, and quantized variants)---rather than a novel
architecture or a new pretraining run. We document the inherited Qwen3-8B architecture faithfully,
describe the fine-tuning pipeline we actually apply, and direct readers to the upstream Qwen3
technical report for the underlying pretraining methodology and benchmark numbers. This report is
intended to give an honest account of what Zen is: a thin, well-documented finetune of a strong
open base model.
\end{abstract}

\tableofcontents
\newpage

%% ─────────────────────────────────────────────────────────────────────────────
\section{Introduction}

Foundation language models trained at scale on diverse internet-scale corpora have demonstrated
remarkable generalization across tasks without task-specific fine-tuning \cite{brown2020gpt3,
wei2022emergent}. Building a competitive base model from scratch, however, requires pretraining
compute that is out of reach for most teams. The pragmatic alternative---adopted here---is to start
from a strong, permissively licensed open base model and invest effort in post-training, packaging,
and deployment.

\paragraph{Provenance.} Zen is a derivative of \textbf{Qwen3-8B}, an 8.2B-parameter dense
decoder-only transformer released by the Qwen team at Alibaba Cloud under the Apache-2.0 license
\cite{qwen3report, qwen3hf}. All architectural choices (layer count, attention configuration,
tokenizer, positional encoding) are inherited unchanged from Qwen3-8B. Earlier versions of this
report described a bespoke ``Zen MoDE'' architecture and an independent 3-trillion-token pretraining
run; those claims were inaccurate and have been removed. Zen does not introduce a new architecture
and was not pretrained from scratch by the Zen LM team.

\paragraph{What Zen adds.} Our actual contributions are:

\begin{itemize}
  \item A supervised fine-tuning (SFT) pass on curated multi-turn instruction--response data, on top
        of the released Qwen3-8B weights.
  \item Optional preference optimization (DPO-style) on a smaller preference set to sharpen
        instruction adherence and refusal behavior.
  \item Packaged, ready-to-deploy artifacts: BF16 SafeTensors, GGUF (llama.cpp) quantizations, and
        MLX builds for Apple Silicon, with documented deployment recipes.
\end{itemize}

Zen occupies the entry tier of the Zen family. Other Zen language models are likewise finetunes of
open Qwen3 base models at different parameter scales (see the Zen family overview).

%% ─────────────────────────────────────────────────────────────────────────────
\section{Architecture (Inherited from Qwen3-8B)}

\subsection{Overview}

Zen inherits the Qwen3-8B architecture without modification: a dense decoder-only transformer with
grouped-query attention, SwiGLU feed-forward blocks, RMSNorm, and rotary position embeddings.
Table~\ref{tab:arch} reproduces the Qwen3-8B hyperparameters as published in the upstream model
configuration \cite{qwen3hf}.

\begin{table}[H]
\centering
\caption{Architecture hyperparameters, inherited unchanged from Qwen3-8B \cite{qwen3hf}.}
\label{tab:arch}
\begin{tabular}{lc}
\toprule
\textbf{Hyperparameter} & \textbf{Value} \\
\midrule
Parameters (total)         & 8.2B \\
Non-embedding parameters   & 6.95B \\
Layers                     & 36 \\
Attention heads (query)    & 32 \\
KV heads (GQA)             & 8 \\
Head dimension             & 128 \\
Hidden dimension           & 4096 \\
FFN intermediate dimension & 12{,}288 \\
Vocabulary size            & 151{,}936 \\
Context length (native)    & 40{,}960 \\
Context length (with YaRN) & 131{,}072 \\
Position encoding          & RoPE ($\theta = 1{,}000{,}000$) \\
Activation function        & SiLU (SwiGLU) \\
Normalization              & RMSNorm \\
Tied embeddings            & No \\
\bottomrule
\end{tabular}
\end{table}

\subsection{Grouped Query Attention}

Qwen3-8B employs Grouped Query Attention (GQA) \cite{ainslie2023gqa} with 32 query heads and 8
key-value heads (head dimension 128). This reduces the KV cache memory footprint by 4$\times$
relative to multi-head attention at equivalent query capacity. Attention is computed as:

\begin{equation}
  \text{Attention}(Q, K, V) = \text{softmax}\!\left(\frac{QK^\top}{\sqrt{d_k}}\right)V
\end{equation}

where each group of query heads shares a single key and value projection.

\subsection{Rotary Position Embeddings}

Rotary position embeddings (RoPE) \cite{su2021rope} encode position as a rotation applied to query
and key vectors. Qwen3-8B uses a RoPE base frequency of $\theta = 1{,}000{,}000$ and a native
maximum position of 40{,}960 tokens, extensible to 131{,}072 tokens via YaRN scaling
\cite{peng2023yarn} as documented upstream \cite{qwen3report}.

\subsection{Feed-Forward Network}

Each transformer layer contains a SwiGLU \cite{shazeer2020glu} feed-forward block:

\begin{equation}
  \text{FFN}(x) = \left(\text{SiLU}(xW_{\text{gate}}) \odot xW_{\text{up}}\right) W_{\text{down}}
\end{equation}

with an intermediate dimension of 12{,}288.

\subsection{Tokenization}

Zen uses the Qwen3 byte-pair encoding (BPE) tokenizer unchanged, with a vocabulary of 151{,}936
tokens and the standard Qwen chat template (system, user, and assistant turn markers). We do not
retrain or extend the tokenizer.

%% ─────────────────────────────────────────────────────────────────────────────
\section{Post-Training Methodology}

This section describes what the Zen LM team actually does on top of the released Qwen3-8B weights.
We do \emph{not} perform large-scale pretraining; for the pretraining corpus, scale, and procedure,
see the upstream Qwen3 technical report \cite{qwen3report}.

\subsection{Supervised Fine-Tuning}

Starting from Qwen3-8B, we apply supervised fine-tuning on a curated mixture of multi-turn
instruction--response data spanning:

\begin{itemize}
  \item General question answering and open-ended generation
  \item Multi-turn dialogue with persona consistency
  \item Code generation and debugging
  \item Summarization and document analysis
  \item Step-by-step reasoning traces
\end{itemize}

SFT trains with a low learning rate (on the order of $2 \times 10^{-5}$) for a small number of
epochs, packing sequences and computing loss only on assistant-turn tokens (response masking). The
goal is to adapt formatting and instruction-following behavior, not to alter the base model's
knowledge.

\subsection{Preference Optimization}

We optionally apply Direct Preference Optimization (DPO) \cite{rafailov2023dpo} on a smaller set of
preference pairs to improve instruction adherence and refusal calibration. DPO avoids the
infrastructure overhead of online RLHF \cite{ouyang2022instructgpt, schulman2017ppo} while still
optimizing a preference objective:

\begin{equation}
  \mathcal{L}_{\text{DPO}} = -\mathbb{E}_{(x, y_w, y_l)}\!\left[\log \sigma\!\left(\beta \log
  \frac{\pi_\theta(y_w \mid x)}{\pi_{\text{ref}}(y_w \mid x)} - \beta \log
  \frac{\pi_\theta(y_l \mid x)}{\pi_{\text{ref}}(y_l \mid x)}\right)\right]
\end{equation}

with the reference policy $\pi_{\text{ref}}$ set to the SFT checkpoint.

%% ─────────────────────────────────────────────────────────────────────────────
\section{Evaluation}

\paragraph{On benchmark numbers.} We do not report fabricated head-to-head benchmark tables. The
capabilities of Zen are, to first order, those of its base model Qwen3-8B; for rigorous,
independently reproducible benchmark results on the Qwen3 base and instruct models---covering MMLU,
GSM8K, MATH, HumanEval, multilingual evaluation, and long-context retrieval---we refer the reader to
the official Qwen3 technical report \cite{qwen3report} and the Qwen3-8B model card \cite{qwen3hf}.
Any task-specific numbers we publish for a given Zen release are measured on the released artifact
under a stated harness and are reported alongside the corresponding Qwen3-8B baseline so that the
effect of our fine-tuning is transparent. We deliberately avoid quoting numbers we have not
measured.

\paragraph{Inheritance, not improvement, by default.} Light SFT and DPO on a strong base model
typically change instruction-following style and formatting more than raw knowledge or reasoning
accuracy. We therefore do not claim that Zen exceeds Qwen3-8B on standard knowledge benchmarks
absent a measured, attributed comparison.

%% ─────────────────────────────────────────────────────────────────────────────
\section{Inference and Deployment}

Because Zen is architecturally identical to Qwen3-8B, it runs on any inference stack that supports
Qwen3 (Hugging Face Transformers, vLLM, llama.cpp via GGUF, and MLX). At 8B parameters the model
fits comfortably on a single 24\,GB consumer GPU (e.g., RTX 4090) in BF16, and runs on commodity
hardware in 4-bit quantization. We publish the following artifacts:

\begin{itemize}
  \item BF16 SafeTensors (reference weights).
  \item GGUF quantizations (Q4\_K\_M, Q8\_0) for llama.cpp on CPU/GPU.
  \item MLX 4-bit builds for Apple Silicon.
\end{itemize}

We do not report fabricated throughput tables here; measured throughput depends on the chosen
runtime, quantization, and hardware, and matches that of any equivalent Qwen3-8B deployment.

%% ─────────────────────────────────────────────────────────────────────────────
\section{Related Work}

Zen sits within the now-standard practice of taking a permissively licensed open base model and
adapting it via instruction tuning. Dense decoder-only transformers have been the dominant paradigm
since GPT-3 \cite{brown2020gpt3}, and the LLaMA series \cite{touvron2023llama} popularized open base
models at the 7--8B scale. The Qwen3 series \cite{qwen3report} that underlies Zen continues this
lineage with GQA, SwiGLU, RoPE, and a large multilingual tokenizer. Instruction tuning via SFT
\cite{wei2022finetuned} and preference optimization via RLHF \cite{ouyang2022instructgpt} and DPO
\cite{rafailov2023dpo} are the standard post-training steps that Zen applies.

%% ─────────────────────────────────────────────────────────────────────────────
\section{Licensing and Attribution}

Qwen3-8B is released by Alibaba Cloud under the Apache-2.0 license, which permits commercial use,
modification, and redistribution subject to attribution and the terms of the license. Zen, as a
derivative, is distributed under the same Apache-2.0 terms with the required upstream attribution and
NOTICE. Users should consult the upstream Qwen3-8B model card and LICENSE for authoritative terms.

%% ─────────────────────────────────────────────────────────────────────────────
\section{Limitations}

Zen inherits the limitations of Qwen3-8B and of autoregressive language models generally. The model
can hallucinate plausible-sounding but incorrect facts, particularly for low-frequency knowledge;
reasoning degrades on long multi-step problems; and alignment is incomplete, so adversarial prompts
can elicit policy-violating outputs. Our light post-training does not remove these limitations, and
in some cases fine-tuning can introduce regressions relative to the base model. Because Zen's
knowledge and core capabilities derive entirely from Qwen3-8B, any limitation of the upstream model
is also a limitation of Zen.

%% ─────────────────────────────────────────────────────────────────────────────
\section{Conclusion}

Zen is an honest, lightly post-trained, well-packaged derivative of the Apache-2.0 Qwen3-8B model.
It does not introduce a new architecture and was not pretrained from scratch; its value lies in
curated instruction tuning and convenient deployment artifacts. We document the inherited Qwen3-8B
architecture transparently and defer to the upstream Qwen3 technical report for pretraining details
and rigorous benchmark numbers. Future Zen releases will continue to build on open Qwen3 base models,
with all provenance and licensing stated explicitly.

%% ─────────────────────────────────────────────────────────────────────────────
\begin{thebibliography}{99}

\bibitem{qwen3report}
Qwen Team, Alibaba Cloud, ``Qwen3 Technical Report,'' \textit{arXiv:2505.09388}, 2025.

\bibitem{qwen3hf}
Qwen Team, ``Qwen3-8B model card and configuration,'' Hugging Face,
\url{https://huggingface.co/Qwen/Qwen3-8B}, 2025.

\bibitem{brown2020gpt3}
T.~Brown et al., ``Language Models are Few-Shot Learners,'' \textit{NeurIPS}, 2020.

\bibitem{wei2022emergent}
J.~Wei et al., ``Emergent Abilities of Large Language Models,'' \textit{TMLR}, 2022.

\bibitem{ainslie2023gqa}
J.~Ainslie et al., ``GQA: Training Generalized Multi-Query Transformer Models,''
\textit{EMNLP}, 2023.

\bibitem{su2021rope}
J.~Su et al., ``RoFormer: Enhanced Transformer with Rotary Position Embedding,''
\textit{arXiv:2104.09864}, 2021.

\bibitem{shazeer2020glu}
N.~Shazeer, ``GLU Variants Improve Transformer,'' \textit{arXiv:2002.05202}, 2020.

\bibitem{peng2023yarn}
B.~Peng et al., ``YaRN: Efficient Context Window Extension of Large Language Models,''
\textit{arXiv:2309.00071}, 2023.

\bibitem{touvron2023llama}
H.~Touvron et al., ``LLaMA: Open and Efficient Foundation Language Models,''
\textit{arXiv:2302.13971}, 2023.

\bibitem{wei2022finetuned}
J.~Wei et al., ``Finetuned Language Models Are Zero-Shot Learners,'' \textit{ICLR}, 2022.

\bibitem{ouyang2022instructgpt}
L.~Ouyang et al., ``Training Language Models to Follow Instructions with Human Feedback,''
\textit{NeurIPS}, 2022.

\bibitem{schulman2017ppo}
J.~Schulman et al., ``Proximal Policy Optimization Algorithms,'' \textit{arXiv:1707.06347}, 2017.

\bibitem{rafailov2023dpo}
R.~Rafailov et al., ``Direct Preference Optimization: Your Language Model is Secretly a Reward
Model,'' \textit{NeurIPS}, 2023.

\end{thebibliography}

\end{document}
